\section{A regularized matrix field}\label{sec:field}

Fix a positive diagonal matrix $K=\operatorname{diag}(k_1,\ldots,k_n)$
from the finite-interval action problem, and put
$\kappa=\max_i k_i$.  We construct a field $M_\delta$ on $\Sym_n$,
the real vector space of symmetric $n\times n$ matrices, and a scalar
primitive $J_\delta$ for its negative.  The difference
$J_\delta(-K)-J_\delta(K)$ will give the boundary contribution in the
action argument.  The other required property is a lower bound for
$\tr(M_\delta(B)(K^2-B^2))$ whenever $M_\delta(B)$ is positive
semidefinite.  The trace comparison of Theorem~\ref{thm:trace} will
supply this bound.

The construction of the field and its primitive, together with the scalar
directional estimates, follows \cite[Section~3]{ScalarCompanion}.  We use a resolvent
normalization that identifies the field directly with the integrated
matrix in Theorem~\ref{thm:trace}; this gives the comparison for positive
values of arbitrary rank in Proposition~\ref{prop:field-comparison}.

Throughout this section $\delta>0$ is fixed.  We use the parameters
$0<\sigma<1$, $p=1+\sigma$, $q=1+1/\sigma$, and
$\beta=\sigma-1/2$, and the normalization $c_\sigma$ from
\eqref{eq:trace-normalization}.

\subsection{Construction and regularity}

For $y\geq0$, define
\begin{equation}\label{eq:field-f}
 f_\delta(y)=\frac{c_\sigma}{\pi}\int_0^\infty t^\beta
 \left(\frac1{t+\delta}-\frac1{t+\delta+y}\right)\,dt.
\end{equation}
For $y<0$, set $f_\delta(y)=f_\delta'(0)y$, the tangent-line
extension at zero.  Differentiation under the integral and the change
of variable $t=(y+\delta)u$ give
\begin{equation}\label{eq:field-f-derivative}
 f_\delta'(y)=a_\sigma(y+\delta)^{\beta-1},\qquad
 a_\sigma=\frac{c_\sigma}{\pi}
       \int_0^\infty\frac{u^\beta}{(1+u)^2}\,du>0
       \quad (y\geq0).
\end{equation}
The integrals converge because $-1/2<\beta<1/2$.
Thus $f_\delta$ is $C^1$, strictly increasing, and concave on $\R$.
Its extension permits us to apply it to every real symmetric matrix,
including matrices with negative eigenvalues.

For $s\in\R$ and $B\in\Sym_n$, let
\begin{equation}\label{eq:field-Q}
 Q_s(B)=s^2I-2sB+K^2.
\end{equation}
We define the field by a symmetric improper integral:
\begin{equation}\label{eq:field-M}
 M_\delta(B)=\lim_{r\to\infty}\int_{-r}^r
       \bigl(f_\delta(Q_s(B))-f_\delta(s^2)I\bigr)\,ds.
\end{equation}
The symmetric cutoff cancels the leading term that is odd in $s$.
The next lemma makes this cancellation quantitative and provides the
primitive used in the action identity.

\begin{lemma}\label{lem:field-regularity}
The limit in \eqref{eq:field-M} exists locally uniformly on $\Sym_n$,
and $M_\delta:\Sym_n\to\Sym_n$ is locally Lipschitz.  There is a
$C^1$ function $J_\delta:\Sym_n\to\R$, which we normalize by
$J_\delta(0)=0$, such that
\begin{equation}\label{eq:field-primitive}
 dJ_\delta(B)[H]=-\tr(M_\delta(B)H)
       \qquad (B,H\in\Sym_n).
\end{equation}
\end{lemma}

\begin{proof}
We first recall a useful estimate for the Hilbert--Schmidt norm.
If $h$ is scalar Lipschitz on an interval containing the spectra of
Hermitian matrices $X$ and $Y$, then
\begin{equation}\label{eq:field-HS-calculus}
 \|h(X)-h(Y)\|_\HS\leq\operatorname{Lip}(h)\|X-Y\|_\HS.
\end{equation}
Indeed, in eigenbases $(e_i)$ of $X$ and $(v_j)$ of $Y$, the
$(i,j)$ entries of these two differences are respectively
$(h(\lambda_i)-h(\nu_j))\langle e_i,v_j\rangle$ and
$(\lambda_i-\nu_j)\langle e_i,v_j\rangle$.  The scalar Lipschitz
bound, followed by summation of their squared absolute values, proves
\eqref{eq:field-HS-calculus}.

Restrict $B$ to a bounded set.  For sufficiently large $|s|$, the
spectrum of $E_s(B)=-2sB+K^2$ lies in $[-a|s|,a|s|]$, with a
constant $a$ independent of $B$.  On this interval put
\[
 h_s(y)=f_\delta(s^2+y)-f_\delta(s^2)-f_\delta'(s^2)y.
\]
By \eqref{eq:field-f-derivative},
\[
 \sup_{|y|\leq a|s|}|h_s'(y)|=O(|s|^{2\beta-3}).
\]
Here the arguments $s^2+y$ are positive for large $|s|$, so the
derivative formula applies.  Since $h_s(0)=0$ and
$\|E_s(B)\|_\HS=O(|s|)$, \eqref{eq:field-HS-calculus} gives
\begin{equation}\label{eq:field-tail}
 f_\delta(Q_s(B))-f_\delta(s^2)I
 =-2s f_\delta'(s^2)B+T_s(B),
\end{equation}
where, uniformly on the chosen bounded set,
\[
 \|T_s(B)\|_\HS=O(|s|^{2\beta-2}),\qquad
 \|T_s(B_1)-T_s(B_2)\|_\HS
 \leq C|s|^{2\beta-2}\|B_1-B_2\|_\HS.
\]
For the first bound, include the term $f_\delta'(s^2)K^2$ in
$T_s$; for the second, use
$E_s(B_1)-E_s(B_2)=-2s(B_1-B_2)$.
The first term in \eqref{eq:field-tail} integrates to zero on
$[-r,r]$, while $2\beta-2=2\sigma-3<-1$ makes both remainder
bounds integrable at infinity.  On bounded $s$-intervals,
\eqref{eq:field-HS-calculus} and the scalar Lipschitz bound on
$f_\delta$ give local Lipschitz continuity directly.  This proves
the convergence and regularity assertions.

To construct the primitive, choose $g_\delta$ with
$g_\delta'=f_\delta$.  The trace differentiation formula is
\[
 d\bigl(\tr g_\delta(X)\bigr)[H]
       =\tr(f_\delta(X)H).
\]
It follows first for polynomials by cyclicity of the trace and then
for $g_\delta$ by polynomial approximation in $C^1$ on compact
intervals containing the relevant spectra.  For $s\ne0$, the function
\[
 B\longmapsto
 \frac{\tr\bigl(g_\delta(Q_s(B))-g_\delta(K^2)\bigr)}{2s}
      +f_\delta(s^2)\tr B
\]
therefore has differential
$-\tr((f_\delta(Q_s(B))-f_\delta(s^2)I)H)$.
At $s=0$ the field is constant and has a linear primitive.
It follows that the integral of the negative field around every
piecewise smooth closed curve is zero at each finite cutoff in
\eqref{eq:field-M}.  Local uniform convergence passes this identity
to $M_\delta$.  Path integration on the vector space $\Sym_n$ now
defines $J_\delta$, with \eqref{eq:field-primitive}; continuity of
$M_\delta$ gives $J_\delta\in C^1$.
\end{proof}

\subsection{Scalar shifts, matrix directions, and endpoints}

The scalar function associated with the field is
\begin{equation}\label{eq:field-mu}
 \mu_\delta(z)=\int_\R
       \bigl(f_\delta(s^2+z)-f_\delta(s^2)\bigr)\,ds,
       \qquad z\in\R.
\end{equation}
This integral is absolutely convergent: on bounded sets of $z$,
its integrand and its Lipschitz constant in $z$ are
$O(|s|^{2\beta-2})$ at infinity.  Consequently $\mu_\delta$ is
locally Lipschitz.  Strict increase of $f_\delta$ implies that
$\mu_\delta$ is strictly increasing, and $\mu_\delta(0)=0$.
For $z\geq0$, Tonelli's theorem and
$\pi^{-1}\int_\R(s^2+a)^{-1}\,ds=a^{-1/2}$ give
\begin{align}
 \mu_\delta(z)
 &=c_\sigma\int_0^\infty t^\beta
       \bigl((t+\delta)^{-1/2}-(t+\delta+z)^{-1/2}\bigr)\,dt
       \notag\\
 &=(z+\delta)^\sigma-\delta^\sigma.
       \label{eq:field-mu-positive}
\end{align}
The last equality is the difference of
\eqref{eq:trace-scaling} at $y=\delta+z$ and $y=\delta$.

We will also translate $s$ inside the first term of
\eqref{eq:field-mu}.  To justify this when the individual terms
are not integrable, fix $z\in\R$ and set
$F_z(s)=f_\delta(s^2+z)$.  This function is even, and
\[
 F_z'(s)=O(|s|^{2\beta-1})\longrightarrow0
             \quad\text{as }|s|\longrightarrow\infty.
\]
For a fixed $a\in\R$ and $r>|a|$, evenness gives the bound
\[
 \left|\int_{-r}^r\bigl(F_z(s-a)-F_z(s)\bigr)\,ds\right|
 \leq |a|^2
       \sup_{r-|a|\leq u\leq r+|a|}|F_z'(u)|.
\]
For example, when $a\geq0$, the difference consists of the integral
over $[r,r+a]$ minus the integral over $[r-a,r]$; pairing these
intervals proves the estimate.  Its right-hand side tends to zero.
Thus
\begin{equation}\label{eq:field-shift}
 \lim_{r\to\infty}\int_{-r}^r
  \bigl(f_\delta((s-a)^2+z)-f_\delta(s^2)\bigr)\,ds
       =\mu_\delta(z).
\end{equation}

These scalar identities control the field in each matrix direction.
They also identify its primitive on the diagonal matrices, where the
two endpoints $K$ and $-K$ lie.

\begin{lemma}\label{lem:field-properties}
For every $B\in\Sym_n$ and every real unit vector $e$,
\begin{equation}\label{eq:field-jensen}
 e^{\mathsf T}M_\delta(B)e
 \leq\mu_\delta\bigl(e^{\mathsf T}K^2e
                         -(e^{\mathsf T}Be)^2\bigr).
\end{equation}
For a standard basis vector $e_i$ and $y\in\R$,
\begin{equation}\label{eq:field-coordinate}
 Be_i=ye_i\quad\Longrightarrow\quad
 M_\delta(B)e_i=\mu_\delta(k_i^2-y^2)e_i.
\end{equation}
If $S$ is a diagonal matrix with diagonal entries in $\{1,-1\}$,
then
\begin{equation}\label{eq:field-equivariance}
 M_\delta(SBS)=S M_\delta(B)S.
\end{equation}
Finally,
\begin{equation}\label{eq:field-endpoints}
 J_\delta(-K)-J_\delta(K)
   =\sum_{i=1}^n\int_{-k_i}^{k_i}
                    \mu_\delta(k_i^2-s^2)\,ds.
\end{equation}
\end{lemma}

\begin{proof}
Scalar concavity, applied in an orthonormal eigenbasis of $Q_s(B)$,
gives
\[
 e^{\mathsf T}f_\delta(Q_s(B))e
      \leq f_\delta(e^{\mathsf T}Q_s(B)e).
\]
Set $a=e^{\mathsf T}Be$ and
$z=e^{\mathsf T}K^2e-a^2$.  Then
$e^{\mathsf T}Q_s(B)e=(s-a)^2+z$.
Subtract $f_\delta(s^2)$, integrate over $[-r,r]$, and use
\eqref{eq:field-shift} to obtain \eqref{eq:field-jensen}.

If $Be_i=ye_i$, the vector $e_i$ is also an eigenvector of $Q_s(B)$,
with eigenvalue $(s-y)^2+k_i^2-y^2$.  Functional calculus and
\eqref{eq:field-shift} prove \eqref{eq:field-coordinate}.
For \eqref{eq:field-equivariance}, use $SK^2S=K^2$ to obtain
$Q_s(SBS)=SQ_s(B)S$, and then conjugate inside
\eqref{eq:field-M}.

On the diagonal subspace, \eqref{eq:field-coordinate} and
\eqref{eq:field-primitive} give
\[
 dJ_\delta(\operatorname{diag}(b_1,\ldots,b_n))
       =-\sum_{i=1}^n\mu_\delta(k_i^2-b_i^2)\,db_i.
\]
Integrating from $(k_1,\ldots,k_n)$ to $(-k_1,\ldots,-k_n)$
within this subspace proves \eqref{eq:field-endpoints}.
\end{proof}

\subsection{The trace estimate at positive values of the field}

We have constructed a locally Lipschitz field and computed the
boundary difference of its primitive.  It remains to connect this
field with Theorem~\ref{thm:trace}.  Positivity of $M_\delta(B)$
will imply positivity of every $Q_s(B)$, making the resolvent
representation available even though $K^2-B^2$ can be indefinite.

\begin{proposition}\label{prop:field-comparison}
For every $B\in\Sym_n$ such that $M_\delta(B)\geq0$,
\begin{equation}\label{eq:field-error}
 \tr\bigl(M_\delta(B)(K^2-B^2)\bigr)
       \geq\tr\bigl(M_\delta(B)^q\bigr)-E_\delta,
 \qquad
 E_\delta=n\bigl(\delta\kappa^{2\sigma}
                 +\delta^\sigma(\kappa^2+\delta)\bigr).
\end{equation}
In particular, $E_\delta\to0$ as $\delta\downarrow0$ for fixed $K$.
\end{proposition}

\begin{proof}
Write $M=M_\delta(B)$.  By \eqref{eq:field-jensen}, its positivity
and the strict increase of $\mu_\delta$ imply
\[
 e^{\mathsf T}K^2e\geq(e^{\mathsf T}Be)^2
                 \qquad\text{for every real unit vector }e.
\]
Consequently
\[
 e^{\mathsf T}Q_s(B)e
  =(s-e^{\mathsf T}Be)^2
       +e^{\mathsf T}K^2e-(e^{\mathsf T}Be)^2\geq0,
\]
and hence $Q_s(B)\geq0$ for all $s\in\R$.  This is positivity
also on $\C^n$: for $z=u+iv$ with real $u,v$, symmetry gives
$z^*Q_s(B)z=u^{\mathsf T}Q_s(B)u+v^{\mathsf T}Q_s(B)v$.
The same directional estimate, together with
\eqref{eq:field-mu-positive}, gives
\begin{equation}\label{eq:field-M-bound}
 0\leq M\leq\mu_\delta(\kappa^2)I
          \leq\kappa^{2\sigma}I.
\end{equation}
The last inequality uses $(a+b)^\sigma\leq a^\sigma+b^\sigma$
for $a,b\geq0$ and $0<\sigma<1$.

Put
\[
 D=K^2-B^2,\qquad V=D+\delta I,\qquad Y_s=(sI-B)^2.
\]
Then $X_s=Y_s+V=Q_s(B)+\delta I\geq\delta I$ satisfies the
hypothesis of Theorem~\ref{thm:trace}.  Let $R(t)$ and $N$ be its
resolvent average and integrated matrix, as defined in
\eqref{eq:trace-resolvents} and \eqref{eq:trace-N}.  We claim that
\begin{equation}\label{eq:field-N}
 N=M+\delta^\sigma I.
\end{equation}

Diagonalize $B$ and apply \eqref{eq:field-shift} with $z=0$ to
each diagonal entry.  This allows $f_\delta(s^2)I$ in
\eqref{eq:field-M} to be replaced by $f_\delta(Y_s)$.
Both $Q_s(B)$ and $Y_s$ are positive semidefinite, so
\eqref{eq:field-f} expresses the resulting difference as
\[
 \frac{c_\sigma}{\pi}\int_0^\infty t^\beta
  \left((Y_s+(t+\delta)I)^{-1}
       -(Q_s(B)+(t+\delta)I)^{-1}\right)\,dt.
\]
For fixed $B$ and $\delta>0$, positivity and quadratic growth in
$s$ give a common lower bound $c(1+t+s^2)I$ for the two matrices
being inverted.  Their difference before inversion is the constant
matrix $D$.  The resolvent identity therefore bounds the norm of
the weighted difference by
\[
 C\|D\|_{\mathrm{op}}\,t^\beta(1+t+s^2)^{-2}.
\]
This is integrable on $(0,\infty)\times\R$: integration in $s$
leaves a constant times $t^\beta(1+t)^{-3/2}$, integrable because
$-1<\beta<1/2$.  Fubini's theorem now yields
\[
 M=c_\sigma\int_0^\infty t^\beta
          \bigl((t+\delta)^{-1/2}I-R(t)\bigr)\,dt.
\]
Subtract this identity from the definition of $N$ and apply
\eqref{eq:trace-scaling} with $y=\delta$.  This proves
\eqref{eq:field-N}.

In particular $N\geq0$, so Theorem~\ref{thm:trace} applies and gives
\begin{align*}
 \tr(MD)
 &=\tr(NV)-\delta\tr M-\delta^\sigma\tr V\\
 &\geq\tr(N^q)-\delta\tr M-\delta^\sigma\tr V.
\end{align*}
Because $N=M+\delta^\sigma I$, the matrices $N$ and $M$ have a
common eigenbasis and $\tr(N^q)\geq\tr(M^q)$.
Equation~\eqref{eq:field-M-bound} bounds $\tr M$ by
$n\kappa^{2\sigma}$, while
$V=K^2-B^2+\delta I\leq(\kappa^2+\delta)I$ bounds $\tr V$ by
$n(\kappa^2+\delta)$.  Substitution proves
\eqref{eq:field-error}.
\end{proof}

In particular, if a matrix $A\in\R^{n\times d}$ satisfies
$M_\delta(B)=AA^{\mathsf T}$, then
Proposition~\ref{prop:field-comparison} controls the nonlinear term
$\tr((AA^{\mathsf T})^q)$ in the action.  The next section constructs
the connecting paths.  In Section~\ref{sec:action-limit}, their action
identities penalize the squared residual
$\|M_\delta(B)-AA^{\mathsf T}\|_\HS^2$ by a factor tending to infinity.
Compactness of pairs of matrices then reduces the estimate to the exact
constraint.
